\section{围绕语义边界组织系统}
\label{sec:design}

系统由服务端、账户守护进程与终端客户端构成。三者的划分并非按界面层次机械拆分，而是对应三种寿命：服务端管理权威事实，守护进程维持账户连接与可恢复的本地视图，终端只管理一次交互会话。由此，关闭窗口不再意味着停止接收消息，恢复界面也不必重新解释网络协议。设计的核心问题随之收敛为：哪些状态必须串行，哪些工作可以并行，以及何时能够向用户宣告一次操作已经成立。

\subsection{以所有权约束并发，而非以线程划分连接}

服务端由单个反应器独占套接字、增量解码器与收发队列，将 SQL 和领域操作交给全局有界工作池。工作线程只返回操作结果，不直接写入连接。该结构借鉴分阶段事件驱动架构对显式队列与过载控制的强调\cite{seda}，但不为每类操作建立独立执行级：对这一规模的聊天负载，单一工作池已经能够隔离数据库阻塞，进一步分级反而增加调度与状态迁移。

每个会话最多提交一个在途任务；后续命令保留在该会话的先进先出队列中。于是，认证、加入房间、发送与读取天然保持客户端提交顺序，而不同连接仍可并行执行。Linux 使用水平触发的 \texttt{epoll}，其余平台使用 \texttt{poll}/\texttt{WSAPoll}；选择水平触发是为了避免把边沿重置与必须彻底排空输入的约束传播到业务层\cite{epoll}。一次就绪事件的处理量受到限制，慢接收者则由独立输出预算约束。

所有权也决定退出语义。收到 TCP FIN 仅表示输入结束：已经完整接收的命令继续执行，回复排空后再关闭套接字。否则，“发送后关闭写端”会被错误解释为取消已提交工作。工作完成通过原生非阻塞套接字唤醒反应器，而不是等待周期轮询；通知缓冲区已满意味着已有待读通知，无须增加重试线程。广播则只编码、压缩一次，各接收者共享不可变字节缓冲区并保有自己的写入偏移。资源核算仍按每个接收者的逻辑字节计费，避免共享存储意外放宽慢客户端限制。

\subsection{提交与发布是同一个顺序问题}

数据库事务解决持久化，却不会自动保证网络观察顺序。若较小编号的消息已经提交、但其工作线程尚未发布结果，另一个线程就可能先广播较大编号。客户端据此推进历史或已读游标，便可能跨过尚未看到的消息。这里缺失的不是另一条 SQL 约束，而是从事务执行到结果入队之间的连续顺序。

系统以租户和房间为键映射到固定的 256 个互斥条带，形成条带锁（striped locking）；相关操作从执行开始持锁，直至完成结果进入反应器的有序完成队列。数据库房间行锁约束事务内部的发送与读取；进程内条带则把提交顺序延伸到事件发布顺序。固定条带不随房间数量增长，代价是哈希碰撞可能串行化无关房间。热点房间本来就需要共享顺序，因此首先保留这一简单结构，而非引入可回收锁表及其额外寿命协议。不同房间不提供全局事件顺序。

PostgreSQL 是权威存储\cite{postgres}。发送必须提交后才生成成功确认或广播；租户、用户和发送随机标识组成唯一键，使原内容重试返回原消息编号而不重复广播，改写内容则产生冲突。身份取自认证会话，而非普通命令字段。数据库连接以独占租约交给工作线程，避免并发使用同一连接。连接故障时不自动重放事务，因为“未收到确认”无法区分“未提交”与“已提交但回复丢失”。这一不确定性必须保留到上层，而不能用盲目重试掩盖。

\subsection{把本地状态绑定到账户，而非窗口}

守护进程把远端凭据与本地 IPC 凭据分离，以账户身份文件和操作系统独占锁约束状态目录。同一目录既不能并发写入，也不能被另一账户静默接管。端点文件原子替换，消息日志不记录认证帧；POSIX 私有权限与 Windows 当前用户专属访问控制分别落实平台边界。多个终端的请求编号在转发时重新映射，响应因而不会因编号碰撞交付给错误窗口。

本地 JSONL 日志是有界缓存，而非第二个数据库：它支持离线浏览、容忍损坏记录，并通过压缩替换控制体积，但不承担提交持久性。重连采用退避，既有终端继续浏览；补齐历史由前台根据独立的已获取游标完成。房间公告的最新编号、实时收到的最大编号与历史分页进度不能合并，否则实时事件会把分页起点推过尚未拉取的区间。

界面发送采用乐观呈现，再按消息编号或作者与随机标识合并确认、广播和历史。尚未转发的离线失败与已经转发但结果丢失必须区分：前者可明确失败，后者显示交付未知，不自动重发。这样的状态划分使用户既能立即看到输入结果，也不会把本地回显误认作数据库承诺。

\subsection{已读是一项界面语义，而非接收副作用}

守护进程收到消息、缓存重放以及历史查询均不推进已读状态。只有前台显式操作才能提交游标，且游标来自当前布局中确实进入视口的消息；历史同步期间暂缓标记。数据库再保证游标单调且不超过房间已提交最大编号。这里的“已读”准确含义是用户确认已显示的消息前缀，而非系统推断用户阅读了每一个字符。将该语义放在布局之后计算，才能使滚动、换行、帮助遮罩和窗口大小真正参与正确性。

终端层使用固定版本的 FTXUI 管理输入、布局和屏幕生命周期\cite{ftxui}，网络仍使用原生套接字。独立有界 I/O 线程合并通知到界面事件循环，使输入节奏不再依赖网络等待。消息正文先消除不可信控制序列；界面只为视口内容建立显示节点，并以有界换行缓存减少长历史重排。与此同时，脚本用纯文本模式保留原有命令契约，使交互升级不破坏既有用户。

这些选择明确界定了系统的伸展方向：单进程承担实时广播，跨进程共享 PostgreSQL 尚不等于分布式订阅；房间未读统计仍需访问历史，持续增长的历史将推动聚合与保留策略；Unicode 宽度估计也不替代完整字素塑形。当前优先保证账户隔离、消息顺序与交互连续性，后续扩展应维护这些已有语义，而不是仅替换底层组件。
